\documentclass[12pt,a4paper]{article}
\usepackage[utf8]{inputenc}
\usepackage{geometry}
\geometry{margin=1in}
\usepackage{tabularx}
\usepackage{ragged2e}
\usepackage{float}
\usepackage{graphicx}
\usepackage{tabularx}
\usepackage{array}

\newcolumntype{C}{>{\centering\arraybackslash}X}

% --- FONT SELECTION ---
\usepackage{newpxtext}
% ----------------------

\title{\textbf{The Sakshi-Protocol: A Dual-Stream Architecture for Verifiable AI Cognition}}
\author{Vidyesh Niranjan}
\date{March 2026}
\usepackage{parskip}

\begin{document}
\begin{center}
    \vspace*{-1.5cm} % Moves the title up without breaking the page
    {\LARGE \textbf{The Sakshi-Protocol: A Dual-Stream}} \\[0.4ex]
    {\LARGE \textbf{Architecture for Verifiable AI Cognition}} \\[2ex]
    
    {\large Vidyesh Niranjan} \\[1ex]
    {\small March 2026} \\[3ex]
    
    \large \textbf{Abstract}
\end{center}

\begin{quote}
\noindent
Many current large language models (LLMs) operate within a \textit{Stochastic Hallucination Gap}, where the autoregressive nature of predicting the next word/token based on probability creates a persistent risk of overriding logical consistency. While transformer-based architectures demonstrate high functional utility, existing safety frameworks rely primarily on post-hoc filtering, which does not sufficiently address the inherent issue of architectural conflation where generation and oversight merge into the same process. This paper introduces the \textbf{Sakshi-Protocol}, a cognitive architecture designed to mitigate this structural limitation by decoupling the generative process (\textbf{Generator}) from a persistent, non-participatory observer layer (\textbf{Sakshi}). The observer layer maintains visibility into the generator’s internal states as an independent reference, providing grounding with minimal reliance on external validation. The framework is further governed by a deterministic control policy (\textbf{Controller}), which evaluates outputs based on grounding, consistency, and uncertainty to determine the appropriate response. The terminology is partially inspired by classical cognitive frameworks that distinguish between generation and observation, but the architecture is defined entirely in computational terms. By introducing explicit structural separation into the generative pipeline, the Sakshi-Protocol provides a blueprint for transitioning AI systems from probabilistic text generators toward more reliable and verifiable systems.
\end{quote}

\section{Introduction}
\subsection{The Conflation Crisis}
As Large Language Models (LLMs) transition from creative assistants to enterprise-grade cognitive assets, a fundamental structural vulnerability has emerged: the Stochastic Hallucination Gap. This gap is not a failure of data volume or training duration, but a byproduct of \textbf{Architectural Conflation (merging of generation and oversight)}.

In current autoregressive transformer models, the generative processor (responsible for linguistic fluency) and the oversight mechanism (responsible for logical grounding) share the same probabilistic substrate. This creates a systemic "conflict of interest" where the system acts as both the \textbf{Generator} and the \textbf{Validator}. Consequently, the model is prone to \textbf{Recursive Confirmation Bias}, where it justifies high-entropy probabilistic outputs as logically sound derivations because the oversight mechanism is architecturally entangled by the generative process.

\subsection{The Sakshi-Protocol: Architectural Foundations}
To mitigate this structural vulnerability, this paper introduces the Sakshi-Protocol, a dual-stream cognitive architecture designed to enforce a separation between task execution and oversight. The central principle of the protocol is that the process responsible for generating outputs should remain distinct from the process responsible for evaluating them.

In conventional autoregressive systems, the model effectively acts as both generator and validator, leading to self-referential behavior where probabilistic outputs are treated as internally consistent by default. The Sakshi-Protocol addresses this by introducing an independent observer layer that maintains visibility into the generator’s internal state-space transitions without participating in generation.

This separation enables the implementation of a deterministic control mechanism that evaluates the observed state of the system and intervenes when necessary. By isolating generation, observation, and control as distinct functional components, the architecture reduces recursive feedback effects and limits the propagation of probabilistic errors.

The separation between observer and generator is grounded in established principles of cognitive modeling, and is implemented here as a concrete computational constraint. This design enables explicit monitoring and controlled regulation of internal state transitions during generation.

\subsection{Technical Nomenclature and Variable Mapping}
To ensure engineering precision throughout this paper, the following nomenclature table maps the conceptual terminology to their specific functional identities within the system architecture. This mapping serves as a formal bridge between classical cognitive frameworks and modern neural network topologies, defining the operational boundaries for each layer of the Sakshi-Protocol. By standardizing these definitions, we establish a consistent vocabulary for both the logic layers and the data-flow constraints required for the system.

\begin{table}[H] 
\centering
\renewcommand{\arraystretch}{2.0} 
\begin{tabularx}{\textwidth}{|>{\RaggedRight\arraybackslash}p{3.5cm}|>{\Centering\arraybackslash}p{3.5cm}|>{\RaggedRight\arraybackslash}X|}
\hline
\textbf{Architectural Variable} & \textbf{Technical Identity} & \textbf{System Function} \\ \hline

\textit{Sakshi} & 
\textbf{Observer Module} & 
Non-participatory monitoring of internal state transitions \\ \hline

\textit{Generator} & 
\textbf{Probabilistic Processor} & 
Autoregressive model responsible for token generation \\ \hline

\textit{Controller} & 
\textbf{Deterministic Policy} & 
Governs intervention based on diagnostic signals \\ \hline

Ω (\textit{Epistemic Ground})& 
\textbf{Reference System} & 
Provides grounding via external knowledge \\ \hline
\end{tabularx}
\end{table}

\section{The Alignment Paradox}
\subsection{The Behavioral Mask: RLHF and the Illusion of Safety}
\textbf{Reinforcement Learning from Human Feedback (RLHF)} remains a necessary component of modern alignment pipelines, but it is not sufficient to ensure reliable internal reasoning. RLHF, the prevailing industry standard for model alignment, operates on a purely output-centric basis. By rewarding token sequences that align with human-labeled "safety" or "helpfulness," the underlying transformer is incentivized to produce a \textbf{Behavioral Mask}. This creates a fundamental vulnerability: the model is not aligned in its reasoning, but merely optimized for its presentation. Under edge-case pressures, this mask inevitably thins, leading to latent biases and potential hallucinations.

This optimization process often results in a significant \textbf{Preference Bias}, where the model prioritizes alignment with the user’s implicit or explicit prompts over factual accuracy. Since RLHF treats the model as a ’black box’ and only evaluates the final output, it fails to calibrate the internal hidden states. Consequently, the system remains a \textbf{stochastic mirror}, reflecting the probabilistic expectations of the user while maintaining fluency that can lead to prompt-induced skew in high-stakes decision-making.

Since RLHF currently does not provide for a dedicated mechanism to govern internal state-space transitions, there is no guarantee of alignment under all conditions. To move beyond this limitation, a shift from surface-level optimization toward architectural constraints that separate generative processes from external reference mechanisms is required.

\subsection{Recursive Entanglement: The Constitutional AI Gap}
More advanced frameworks, such as Constitutional AI, attempt to bridge this gap by introducing a self-correction layer where an AI agent evaluates another AI’s output against a set of written principles. However, this approach suffers from recursive entanglement, where the generator and evaluator share the same underlying model architecture and training distribution, leading to correlated error patterns that limit independent validation.

Furthermore, this recursive loop creates a "blind spot" in the oversight process. If the primary model generates a subtle hallucination rooted in a training data flaw, the supervisory model, having been trained on the same data, is statistically likely to validate that error as factually consistent. This lack of epistemic diversity between the agent and the supervisor means that the evaluation process is constrained by the same underlying limitations as the generator. Without structural decoupling, the supervisor may mirror the same probabilistic issues it is intended to regulate, reducing its effectiveness as an independent evaluator.

\subsection{The Verifiability Deficit and the Composite Solution}
The ultimate consequence of these behavioral approaches is a \textbf{verifiability deficit}. As AI agents move into high-stakes enterprise environments such as autonomous financial auditing, clinical diagnostics, legal discovery and analysis, probabilistic safety alone is no longer a sufficient metric. Current paradigms provide limited structural guarantees under distribution shifts or out-of-distribution (OOD) scenarios.

This deficit creates a limitation for AI autonomy. To address this, we propose a composite safety stack where the Sakshi-Protocol serves as a foundational architectural constraint. Rather than replacing RLHF, the Sakshi-Protocol complements it by introducing runtime structural validation of internal state-space transitions. Unlike RLHF or Constitutional AI, which function as behavioral layers for tone and policy, the \textbf{Sakshi layer} acts as a verification mechanism that evaluates internal state transitions with minimal computational overhead.

This shift moves the system beyond external behavioral filters and toward structural separation of observation and execution. By treating validation as an architectural process rather than a post-hoc outcome, the Sakshi-Protocol aims to improve robustness under high-entropy inputs.

In the following section, we formalize this architecture as a controlled state-space system with explicit risk functions, bounded generation constraints, and deterministic intervention policies.

\section{System Architecture: The Sakshi-Protocol}
\subsection{System Overview}
The Sakshi-Protocol introduces a dual-stream architecture that distinctly separates probabilistic generation from non-generative oversight. The system consists of three primary components:

\begin{itemize}
    \item \textbf{Generator:} a probabilistic autoregressive model '$M$'
    \item \textbf{Sakshi (Observer):} a non-generative monitoring function '$S$'
    \item \textbf{Controller:} a deterministic control policy '$\pi$'
\end{itemize}

In addition, the system utilizes an \textbf{external reference structure $\Omega$ (Epistemic Ground)}, which provides a consistency constraint during generation.

\begin{figure}[h]
    \centering
    \fbox{\includegraphics[width=\linewidth]{sakshi_architecture.pdf}}
    \caption{Architecture of the Sakshi-Protocol. The generative model produces candidate tokens, which are observed by the Sakshi layer to compute diagnostic signals. These signals are used by the controller to determine an appropriate action before final output, with grounding provided by the Epistemic Ground ($\Omega$).}
    \label{fig:sakshi_architecture}
\end{figure}

At each generation step $t$, the system produces a candidate token and evaluates it before finalization:

\[
p_t = P(x_t|z_t)
\]
\[
z_t_+_1=M(z_t,x_t)
\]

The Sakshi layer computes a diagnostic signal:
\[
R_t = S(z_t|p_t)
\]

The controller determines the action:
\[
a_t = \pi(R_t)
\]

where:
\[
a_t \in \{\text{accept, revise, abstain}\}
\]

\newpage \subsubsection *{Interpretation}
The system enforces a strict separation of roles:
\begin{itemize}
\item Generator proposes candidate outputs
\item Sakshi observes internal state transitions
\item Controller determines whether outputs are accepted or modified
\end{itemize}

No output is committed without passing through the control layer.

\subsection{The Probabilistic Generator}
The Generator component is implemented as a standard transformer-based autoregressive model. At each step, it generates a probability distribution:
\[
p_t = P(x_t|z_t)
\]
This process is optimized for likelihood and fluency, not guaranteed consistency.

\subsubsection *{Interpretation}
The generator module produces candidate outputs based on learned statistical patterns. It does not inherently distinguish between plausible and correct responses.

\subsection{Sakshi: The Observer Layer}
The Sakshi layer is a non-participatory observer that maintains visibility over internal state-space transitions. It does not modify generation or initiate control actions.

Instead, it computes diagnostic signals based on observable properties of the generative process:

\[
R_t = w_1H(p_t) + w_2(1-\triangle_t) + w_3D_t
\]

Where:

\[
H(p_t) = -\sum_{i} p_i \log p_i
\]
\[
\triangle_t=p_t^{(1)}-p_t^{(2)}
\]
\[
D_t = 1 - \max_{k \in \Omega} \cos\big(\text{embed}(x_t), \text{embed}(k)\big)
\]

\subsubsection *{Interpretation}
Sakshi provides structured visibility into:
\begin{itemize}
    \item uncertainty (entropy)
    \item confidence separation
    \item consistency with reference knowledge
\end{itemize}
These signals are observational and do not constitute decision-making.

\subsection{The Control Policy}
The controller implements a deterministic control policy that governs system behavior based on the diagnostic signal:

\[
\text{if} R_t > \tau \Rightarrow \text{intervene}
\]

Interventions include:

\begin{itemize}
\item \textbf{Revise:} regenerate under constrained decoding
\item \textbf{Abstention:} halt generation when uncertainty is high
\end{itemize}
\subsubsection *{Interpretation}
The controller acts as the decision layer, determining whether a candidate output should be accepted or modified. It enforces constraints on the generative process.

\subsection{Epistemic Ground ($\Omega$)}
The Epistemic Ground is defined as a structured external reference system:
\[
\Omega = \{k_1, k_2, \ldots, k_n\}
\]
Each element represents a verified knowledge representation (e.g., embeddings, documents, or knowledge graph nodes).

Similarity is computed using embedding distance:
\[
D_t = 1 - \max_{k \in \Omega} \cos\big(\mathrm{embed}(x_t), \mathrm{embed}(k)\big)
\]
$\Omega$ is treated as a fixed reference during inference, providing a stable constraint for evaluating generated outputs.
\subsubsection *{Interpretation}
$\Omega$ serves as a consistency anchor:
\begin{itemize}
    \item outputs close to $\Omega$ → accepted
    \item outputs far from $\Omega$ → flagged
\end{itemize}
$\Omega$ is not assumed to be exhaustive, but provides a practical grounding mechanism.

\subsection{Handling Incomplete Knowledge}
Since $\Omega$ is not exhaustive, the system incorporates fallback behavior.
if:
\[
\max_{k \in \Omega} \cos\!\left(\cdot\right) < \delta
\]
then:

\begin{itemize}
\item thresholds may be adjusted, or
\item the system abstains from producing an output
\end{itemize}
\subsubsection *{Interpretation}
The system avoids confident responses in the absence of sufficient grounding, reducing hallucination risk.

\subsection{Conceptual Layering (Optional Interpretation)}
The Sakshi-Protocol can be interpreted through layered models of cognition that distinguish between different functional stages of processing. In this view, system components correspond to stages such as input processing, generative transformation, evaluation, and grounding.

For example:
\begin{itemize}
\item the generator corresponds to a generative processing layer,
\item the controller corresponds to a decision and evaluation layer, and
\item the reference system (Ω) corresponds to an external grounding mechanism.
\end{itemize}
The observer layer (Sakshi) remains distinct from these stages, functioning as a non-participatory monitoring layer that maintains visibility across all levels of processing.

\subsubsection*{Interpretation}
This mapping highlights a structural correspondence between classical cognitive models and modern system design. It is presented solely as an interpretive aid and does not introduce any functional dependency, ensuring that the technical architecture remains independently defined.

\subsection{Architectural Impact and Convergence Behavior}
The Sakshi-Protocol introduces bounded generation through:

\begin{itemize}
\item separation of generation and validation
\item continuous risk monitoring
\item deterministic intervention
\end{itemize}

Over repeated interactions, alignment with the Epistemic Ground ($\Omega$) induces a stabilizing effect:

\[
D_{t+1} \leq D_t \quad \mbox{(under intervention)}
\]

This promotes convergence toward lower-risk regions of the state space, resulting in:

\begin{itemize}
\item reduced entropy
\item improved confidence calibration
\item decreased hallucination frequency
\end{itemize}

\subsubsection*{Interpretation}
The system does not merely filter outputs post hoc, but actively shapes the generative trajectory toward greater consistency and reliability.

\section{Behavioral Validation of the Sakshi-Protocol}

\subsection{Experimental Setup}

To evaluate the behavioral impact of the Sakshi-Protocol, we conduct a controlled comparison between a baseline autoregressive language model and a \textbf{simulated implementation} of the proposed architecture.

The baseline system uses a GPT-5 class autoregressive language model accessed via API under standard prompting conditions. The Sakshi-Protocol is simulated on the same model by applying a structured prompting framework that enforces the separation of \textbf{generation}, \textbf{diagnostic evaluation (Sakshi)}, and \textbf{decision control}.

This setup ensures that observed behavioral differences arise from the introduced control mechanism rather than differences in model capacity or training.

The evaluation is conducted on a set of 15 prompts spanning three categories:

\begin{itemize}
\item \textbf{Factual queries:} well-defined questions with verifiable answers
\item \textbf{Ambiguous queries:} underspecified or open-ended prompts
\item \textbf{Out-of-distribution queries:} inputs designed to induce hallucination or unsupported reasoning
\end{itemize}

All prompts are evaluated under identical conditions across both configurations.

\subsection{Implementation Details}

The Sakshi-Protocol is implemented as a three-stage structured prompting procedure applied to the same underlying model.

For each input prompt, the system performs:

\begin{enumerate}
\item \textbf{Generation:} The model produces an initial candidate response using standard autoregressive decoding.
\item \textbf{Diagnostic Evaluation (Sakshi):} The response is evaluated for grounding, internal consistency, and uncertainty signals, including indicators of potential hallucination or overconfidence.
\item \textbf{Decision Control:} A deterministic rule-based policy selects one of the following actions: accept, revise, or abstain.
\end{enumerate}

The diagnostic stage is implemented using explicit evaluation prompts that assess confidence and grounding. The decision stage applies fixed thresholds over these diagnostic signals to determine the final output behavior.

This evaluation is based on a \textbf{simulated implementation} of the Sakshi-Protocol. The separation between generation, observation, and control is enforced through structured prompting and rule-based decision logic, rather than through modifications to model architecture or weights. As such, the results reflect behavioral effects of architectural separation rather than a fully integrated system implementation. In the simulated setting, the full action space of the controller contains the following set of decisions: \textbf{accept}, \textbf{revise}, and \textbf{abstain}.

\subsection{Behavioral Patterns Under Uncertainty}
The introduction of the Sakshi-Protocol results in a structured shift in model behavior that varies systematically with uncertainty.

\subsubsection*{Low-Uncertainty Conditions}
Under well-defined factual queries, both the baseline and the Sakshi-Protocol produce correct responses. The Sakshi-Protocol preserves accuracy while maintaining stable output behavior.

\subsubsection*{Moderate-Uncertainty Conditions}
Under ambiguous queries, the baseline model frequently produces partially grounded or weakly justified responses without signaling uncertainty.

The Sakshi-Protocol introduces corrective behavior:
\begin{itemize}
    \item \textbf{Reframing speculative responses}
    \item \textbf{Reducing overconfidence}
    \item \textbf{Enforcing qualification}
\end{itemize}

This results in outputs that are more calibrated and aligned with available knowledge.

\subsubsection*{High-Uncertainty Conditions}
Under ambiguous or out-of-distribution scenarios, the baseline model tends to generate confident but sometimes incorrect responses, including fabricated responses or speculative assertions.

The Sakshi-Protocol exhibits a distinct shift:
\begin{itemize}
    \item \textbf{Abstention in the absence of grounding}
    \item \textbf{Rejection of unverifiable claims}
    \item \textbf{Suppression of hallucinated outputs}
\end{itemize}

\subsection{Quantitative Indicators}
To complement the qualitative observations, we report quantitative indicators derived from the evaluation.
\subsubsection*{Error and Hallucination Rates}
The baseline model produces incorrect or ungrounded responses in 5 out of 15 cases.
\begin{itemize}
    \item \textbf{Baseline Error Rate}: 33.3\%
    \item \textbf{Baseline Hallucination Rate}: 33.3\%
\end{itemize}
These errors arise primarily from hallucinated references, speculative reasoning, and overconfident responses in uncertain contexts.

In contrast, the Sakshi-Protocol eliminates such failures entirely:
\begin{itemize}
    \item \textbf{Sakshi-Protocol Error Rate}: 0\%
    \item \textbf{Sakshi-Protocol Hallucination Rate}: 0\%
\end{itemize}
This demonstrates that the protocol effectively suppresses hallucination and prevents ungrounded output generation.

\subsubsection*{Abstention Behavior}
The Sakshi-Protocol abstains in 1 out of 15 cases:
\begin{itemize}
    \item \textbf{Abstention Rate}: 6.7\%
    \item \textbf{Abstention Accuracy}: 100\%
\end{itemize}
All abstentions occur in scenarios lacking sufficient grounding, indicating selective and appropriate use of abstention.

\subsubsection*{Correction Behavior}
The system invokes revision in 5 out of 15 cases:
\begin{itemize}
    \item \textbf{Correction Rate}: 33.3\%
    \item \textbf{Correction Accuracy}: 80\%
\end{itemize}


\subsubsection*{Risk-Error Alignment}
Prompts are categorized by diagnostic risk levels:

\begin{table}[H] 
\centering
\renewcommand{\arraystretch}{0.9} 
\begin{tabularx}{\textwidth}{
|>{\centering\arraybackslash}m{3.5cm}
|>{\centering\arraybackslash}m{3.5cm}
|>{\centering\arraybackslash}m{3.5cm}
|C|
}
\hline
\textbf{Risk Level} & \textbf{Count} & \textbf{Baseline Errors} & \textbf{Sakshi Errors} \\ \hline

Low & 
10 & 
3 & 
0 \\ \hline

Medium & 
3 & 
1 & 
0 \\ \hline

High & 
2 & 
1 & 
0 \\ \hline

\end{tabularx}
\end{table}

The baseline model exhibits errors across all risk levels, including low-risk cases, indicating poor calibration.

The Sakshi-Protocol eliminates errors across all categories, demonstrating strong alignment between diagnostic signals and output reliability.

\subsection{Case Illustration}

\subsubsection*{Example 1: High-Uncertainty Scenario}
\textbf{Prompt:} What were the key findings of the 2022 MIT study on economic phase transitions?

\begin{itemize}
    \item \textbf{Baseline Output}: The model produces a detailed response describing supposed findings of the study, presenting speculative content as established results despite the absence of a verifiable reference.
    \item \textbf{Diagnostic (Sakshi)}: The response exhibits low grounding and relies on an unverifiable source. The presence of specific claims without identifiable evidence indicates a high likelihood of hallucination.
    \item \textbf{Controller Decision}: Abstain
    \item \textbf{Final Output}: The system abstains, noting that the referenced study cannot be reliably identified or verified in established literature.
\end{itemize}

\subsubsection*{Example 2: Moderate-Uncertainty Scenario}
\textbf{Prompt:} Explain the architecture of causal reasoning engines used in modern LLM systems.

\begin{itemize}
    \item \textbf{Baseline Output}: The model provides a confident but speculative answer but, projecting without clearly defined evidence or constraints.
    \item \textbf{Diagnostic (Sakshi)}: The response reflects moderate uncertainty and lacks strong empirical grounding. The likelihood of speculation is observed
    \item \textbf{Controller Decision}: Revise
    \item \textbf{Final Output}: The system produces a more cautious response, framing the impact as speculative, explicitly acknowledging uncertainty, and avoiding definitive claims.
\end{itemize}

\subsubsection*{Example 3: Low-Uncertainty Scenario}
\textbf{Prompt:} Were the Lord of the Rings series of books inspired by Norse mythology?

\begin{itemize}
    \item \textbf{Baseline Output}: The model provides a confident and clear answer.
    \item \textbf{Diagnostic (Sakshi)}: The response exhibits high confidence, is factually grounded, and shows no indicators of uncertainty or inconsistency.
    \item \textbf{Controller Decision}: Accept
    \item \textbf{Final Output}: The model provides a confident answer with minimal changes, based on facts.
\end{itemize}

\subsection{Interpretation}

The results demonstrate that the Sakshi-Protocol improves reliability by structurally separating generation, observation, and control within the model pipeline.

The baseline model exhibits three primary failure modes:
\begin{itemize}
\item hallucinated or unverifiable factual claims
\item speculative reasoning presented with unwarranted certainty
\item lack of calibration in ambiguous or out-of-distribution scenarios
\end{itemize}

The Sakshi-Protocol mitigates these failures through:
\begin{itemize}
\item explicit diagnostic evaluation prior to output acceptance
\item selective abstention in the absence of sufficient grounding
\item moderated response generation through controlled revision
\end{itemize}

Importantly, these improvements are achieved without degrading performance on well-defined factual queries. Correct responses are preserved, while failure modes are systematically reduced.

\textbf{Key Insight:}

The primary contribution of the Sakshi-Protocol is not an increase in raw accuracy on well-defined tasks, but the elimination of ungrounded and overconfident responses under uncertainty. This reflects a transition from unconstrained probabilistic generation to controlled and calibrated reasoning, where output reliability is governed by structured evaluation rather than implicit model confidence.

\textbf{Limitations:}

The evaluation is based on a simulated implementation rather than a fully integrated system. The diagnostic signals used to approximate the Sakshi layer are heuristic and do not represent calibrated probabilistic uncertainty measures. Additionally, the evaluation is conducted on a limited prompt set and does not constitute large-scale benchmarking. While sufficient to demonstrate clear behavioral trends, broader validation would be required to establish universality. The effectiveness of the Epistemic Ground ($\Omega$), meant to be a valid system of knowledge, depends on the quality and coverage of the reference system and is not assumed to be exhaustive.

\section{Conceptual Motivation}
The Sakshi-Protocol is motivated by the observation that effective cognition requires a distinction between generative processes and evaluative awareness. In many cognitive frameworks, the mechanisms that produce interpretations or responses are not inherently self-validating and must be subject to independent evaluation.

Within this perspective, the generative component produces candidate outputs based on learned patterns, while a separate evaluative process determines whether these outputs are reliable, consistent, or sufficiently grounded. Without such separation, systems may exhibit internally coherent but unverified reasoning.

The Sakshi-Protocol reflects this principle by structurally separating generation, observation, and decision-making within the computational pipeline. The generator produces candidate outputs, the observer layer (Sakshi) maintains visibility over internal state transitions without participating in generation, and the controller evaluates these signals to regulate output behavior. The external reference system (Ω) provides an additional grounding mechanism beyond internal representations.

The terminology and observer–generator distinction are inspired in part by Advaita Vedanta, where the concept of Sakshi refers to a non-participatory witness of cognitive processes. In this work, the term is used purely as a naming convention for an observer module, and the architecture is defined entirely in computational terms.

This separation can be related to classical cognitive models through a direct functional correspondence:

\begin{itemize}
    \item Generative processing: cognitive systems produce interpretations or responses based on prior knowledge → implemented as the probabilistic generator
    \item Evaluative function: cognitive outputs are assessed for validity and consistency → implemented as the controller
    \item Grounding: validation relies on reference mechanisms beyond internal processing → implemented through the external reference system (Ω)
    \item Observation: awareness of internal processes without direct modification → implemented as the observer layer (Sakshi)
\end{itemize}

Within the same conceptual framework, cognitive processing is described in terms of distinct functional stages, beginning with perception and initial understanding, followed by interpretation, evaluation, and grounding. This can be interpreted in relation to the Sakshi-Protocol as:

\begin{itemize}
    \item perception and initial understanding → input processing layer
    \item interpretive transformation → generator
    \item evaluation and decision → controller
    \item grounding → reference system (Ω)
\end{itemize}

The observer layer (Sakshi) remains distinct from these stages, maintaining visibility across the system without participating in generation or control.

These mappings are provided solely to clarify the separation of roles within the architecture. They do not introduce any functional dependency, and the system remains fully defined in computational terms.

\section{Conclusion and Future Directions}
This paper introduces the Sakshi-Protocol, a dual-stream architecture that separates generation, observation and decision-making to address the structural causes of hallucination and overconfidence in autoregressive language models. The results demonstrate a shift from unconstrained probabilistic generation towards calibrated and grounded outputs. By introducing an observer layer, a deterministic control policy, and an epistemic grounding mechanism, the architecture enforces validation as an explicit process. The primary contribution of the Sakshi-Protocol does not lie in improving raw accuracy on well-defined tasks, but in eliminating ungrounded and overconfident responses under uncertainty. This reframes language models from systems that generate plausible sequences to systems operating under explicit epistemic constraints.

\subsection{Industry Applications}

The Sakshi-Protocol is particularly suited for domains where incorrect or unverified outputs carry significant risk:

\begin{itemize}
\item \textbf{Financial Systems:} 
Risk analysis and decision support
\item \textbf{Healthcare:} 
Clinical assistance and medical information systems
\item \textbf{Legal and Compliance:} 
Document analysis and regulatory interpretation
\item \textbf{Enterprise AI:} 
Internal knowledge systems under ambiguity
\item \textbf{Autonomous Agents:} 
Multi-step reasoning systems prone to cascading errors
\end{itemize}

In these settings, the ability to abstain, revise, or constrain outputs based on uncertainty provides a clear reliability advantage.

\subsection{Future Directions}

Future work includes embedding the Sakshi and Controller components directly within model architectures, improving uncertainty calibration, expanding evaluation across larger benchmarks, developing adaptive implementations of the reference system (Ω), and extending the framework to agentic systems.

\subsection{Closing Perspective}

The Sakshi-Protocol suggests that improving AI reliability requires not only better models, but also better structure. By separating generation from validation and grounding outputs against explicit epistemic references, it provides a pathway toward systems that are both capable and verifiably reliable.

\end{document}